% titel: Blatt 7, A1: Maxwell-Rad
% skript: Physik I, Kapitel 6, Seiten 96--100 und 102--110
% stand: Übungsblatt 7, Aufgabe A1

Diese Werkzeuge beschreiben die Bewegung ausgedehnter Körper, bei denen sich
Schwerpunktbewegung und Drehung überlagern. Die Beziehungen gelten zunächst
allgemein. Wähle die passenden Kräfte, Achsen und Vorzeichen selbst.

\begin{werkzeugkasten}[Kräfte und Drehmomente]
Die Bewegung des Schwerpunkts folgt aus der Summe der äusseren Kräfte. Ein
Drehmoment hängt dagegen auch vom Angriffspunkt der Kraft und der gewählten
Drehachse ab:
\begin{equation}
  \sum \vec F_{\mathrm{ext}}=m\vec a_S,
  \qquad
  \vec M_S=\vec r\times\vec F,
  \qquad
  \sum M_S=I_S\alpha.
\end{equation}
Hier bezeichnet $S$ den Schwerpunkt, $I_S$ das Trägheitsmoment um die
Schwerpunktachse und $\alpha$ die Winkelbeschleunigung. Die letzte Beziehung
gilt für die hier relevante Rotation um eine feste Hauptrichtung. Der Betrag des
Drehmoments ist Kraft mal senkrechter Hebelarm. Bei mehreren Kräften müssen
ihre Beiträge addiert werden. Kräfte auf verschiedene Körper nicht in derselben
Kräftebilanz zusammenzählen.

\textit{Skript:} Kap. 6.1 und 6.2.4, S. 96--98 und 102.
\end{werkzeugkasten}

\begin{werkzeugkasten}[Trägheit und Bewegungsenergie]
Für einen homogenen Vollzylinder beziehungsweise eine homogene Scheibe um die
Symmetrieachse gilt $I_S=\tfrac12 mR^2$. Das Trägheitsmoment misst, wie die
Masse bezüglich der Drehachse verteilt ist. Die kinetische Energie eines starren
Körpers setzt sich aus Schwerpunkt- und Rotationsenergie zusammen:
\begin{equation}
  E_{\mathrm{kin}}=\frac12 mv_S^2+\frac12 I_S\omega^2.
\end{equation}
In einem homogenen Schwerefeld ist die potenzielle Energie $mgh_S$. Ohne
Energieverluste und äussere Arbeit bleibt die Summe dieser Beiträge erhalten.
Die Rotationsenergie darf nicht mit der Energie der Schwerpunktbewegung
verwechselt werden.

\textit{Skript:} Kap. 6.2.1--6.2.2 und 6.3, S. 99--101 und 109--110.
\end{werkzeugkasten}

\clearpage
\begin{werkzeugkasten}[Kinematik und Impuls]
Wenn eine Schnur ohne Gleiten an einer zylindrischen Oberfläche mit Radius
$b$ abläuft, sind die \emph{Beträge} von Verschiebung und Drehwinkel durch
$s=b\varphi$ verbunden. Entsprechend gilt $v=b\omega$; während einer Phase
mit unverändertem Kontakt gilt für die Beträge auch $a=b\alpha$. Die
Vorzeichen müssen zur jeweils gewählten Orientierung passen. Der wirksame
Radius ist der Radius am Kontakt, nicht automatisch der äussere Körperradius.

Für eine gleichmässig beschleunigte geradlinige Bewegung gilt
$v^2-v_0^2=2a\Delta s$, mit vorzeichenrichtig gewählter Koordinate. Für einen
Vorgang über ein Zeitintervall $\Delta t$ verknüpft der Impulssatz die
\emph{resultierende} Kraft mit der Impulsänderung:
\begin{equation}
  \int_{t_1}^{t_2}\sum F_{\mathrm{ext}}\,\mathrm dt
    = m(v_2-v_1),
  \qquad
  \overline{\sum F_{\mathrm{ext}}}
    =\frac{m(v_2-v_1)}{\Delta t}.
\end{equation}
Eine mittlere Kraft benötigt also Informationen über die Dauer und die
Geschwindigkeiten des betrachteten Vorgangs. Der Kraftstoss und eine mittlere
Kraft sind unterschiedliche Grössen.

\textit{Skript:} Kap. 6.2.5 und 6.3, S. 106 und 109--110; Impulssatz
Kap. 2.2.2, S. 27--28, und Kap. 3.2, S. 52--53.
\end{werkzeugkasten}

\begin{achtung}[Darauf achten]
\begin{itemize}
  \item Bezeichne die Kraft einer einzelnen Schnur und die Gesamtkraft mehrerer
        Schnüre eindeutig. Kraft und Gegenkraft wirken auf verschiedene Körper.
  \item Bewegung und Beschleunigung müssen nicht in dieselbe Richtung zeigen.
        Für unterschiedliche Bewegungsphasen die Kräfte und den Kontakt neu
        prüfen, statt eine Vorzeichenkonvention ungeprüft zu übernehmen.
  \item Eine ideale, verlustfreie Bewegung ist ein Modell. Ein kurzer
        Umkehrvorgang benötigt gegebenenfalls zusätzliche Angaben; fehlende
        Grössen dürfen nicht durch eine beliebige Zahl ersetzt werden.
\end{itemize}
\end{achtung}
